\subsubsection{Vaizdų tinklas (2D CNN)}
\label{subsubsec:image_cnn}

Klasė \texttt{ImageCNN} naudoja \texttt{nn.Conv2d}, \texttt{nn.BatchNorm2d}, \texttt{nn.MaxPool2d} ir \texttt{nn.Dropout2d} sluoksnius. Įėjimas – $3 \times 64 \times 64$ formato RGB tenzorius (vaizdo kraštinė valdoma konstanta \texttt{IMAGE\_SIZE} $= 64$). Bazinė konfigūracija turi tris konvoliucinius blokus su filtrų skaičiais $(32,\,64,\,128)$, branduolio dydžiu $3 \times 3$ ir telkimo dydžiu $2 \times 2$. Po trijų telkimo sluoksnių požymių žemėlapio kraštinė sumažėja iki $64 / 2^{3} = 8$, todėl po suplojimo gaunamas $128 \cdot 8 \cdot 8 = 8192$ dydžio vektorius, kuris paduodamas į $256$ neuronų paslėptąjį sluoksnį, o galiausiai – į $2$ neuronų išėjimo sluoksnį (klasės \textit{chihuahua} ir \textit{muffin}). Sluoksnių parametrai pateikti \ref{tab:image_cnn_architecture} lentelėje.

\begin{table}[H]
    \centering
    \caption{Vaizdų klasifikavimo tinklo (\texttt{ImageCNN}) bazinės architektūros sluoksniai.}
    \begin{tabular}{|c|l|l|}
        \hline
        Nr. & Sluoksnis & Parametrai / išėjimas \\
        \hline
        $0$ & Įėjimas & $3 \times 64 \times 64$ \\ \hline
        $1$ & \texttt{Conv2d} & $3 \to 32$, branduolys $3$, \texttt{padding}~$=1$ \\ \hline
        $2$ & \texttt{ReLU} & – \\ \hline
        $3$ & \texttt{MaxPool2d} & dydis $2$, išėjimas $32 \times 32 \times 32$ \\ \hline
        $4$ & \texttt{Conv2d} & $32 \to 64$, branduolys $3$, \texttt{padding}~$=1$ \\ \hline
        $5$ & \texttt{ReLU} & – \\ \hline
        $6$ & \texttt{MaxPool2d} & dydis $2$, išėjimas $64 \times 16 \times 16$ \\ \hline
        $7$ & \texttt{Conv2d} & $64 \to 128$, branduolys $3$, \texttt{padding}~$=1$ \\ \hline
        $8$ & \texttt{ReLU} & – \\ \hline
        $9$ & \texttt{MaxPool2d} & dydis $2$, išėjimas $128 \times 8 \times 8$ \\ \hline
        $10$ & \texttt{Flatten} & $8192$ \\ \hline
        $11$ & \texttt{Linear} & $8192 \to 256$ \\ \hline
        $12$ & \texttt{ReLU} & – \\ \hline
        $13$ & \texttt{Dropout} & $p = 0{,}0$ (bazinė) \\ \hline
        $14$ & \texttt{Linear} & $256 \to 2$ \\ \hline
    \end{tabular}
    \label{tab:image_cnn_architecture}
\end{table}

